% Generated draft: requires mathematical and source review.
\documentclass[11pt]{article}
\usepackage[margin=25mm]{geometry}
\usepackage{fontspec}
\setmainfont{Latin Modern Roman}
\usepackage{amsmath,amssymb,mathtools,unicode-math}
\setmathfont{Latin Modern Math}
\usepackage{enumitem,needspace,xurl,hyperref}
\setlength{\parindent}{0pt}
\setlength{\parskip}{5pt}
\setlength{\emergencystretch}{3em}
\newcommand{\N}{\mathbb{N}}
\newcommand{\Z}{\mathbb{Z}}
\newcommand{\Q}{\mathbb{Q}}
\newcommand{\R}{\mathbb{R}}
\newcommand{\C}{\mathbb{C}}
\newcommand{\F}{\mathbb{F}}
\newcommand{\A}{\mathbb{A}}
\newcommand{\PP}{\mathbb{P}}
\newcommand{\E}{\mathbb{E}}
\newcommand{\eps}{\varepsilon}
\newcommand{\dd}{\,\mathrm{d}}
\begin{document}
\section*{UnsolvedMath ID: 20002384}
\textbf{Source problem code:} catalogue entry
\subsection{Definitions and mathematical statement}
\textbf{Participant context, not a posed problem.} In item A.12, Yasuo Yoshinobu says that he works in core model theory and shares the interests in the organizers' detailed conference description. The item itself gives no mathematical hypotheses or requested conclusion. The supplied aid associates that description with Woodin's analysis of HOD in models of \(\mathrm{AD}^{+}\) by extenders and iteration strategies. The detailed announcement and any particular proposed extender/strategy presentation are not reproduced in the record.

\textbf{Supplementary HOD diagnostic reported by the aid.} Let \(M\) be a transitive model of \(\mathrm{ZF}\). Define \(\operatorname{OD}^M\) using ordinal definability as computed internally in \(M\), and let
\[
 \operatorname{HOD}^M=
 \{x\in M:\operatorname{tc}^M(\{x\})\subseteq\operatorname{OD}^M\}.
\]
Let \(\operatorname{Ord}^M\) be its ordinals, and use the real coding
\[
 \mathcal R^M=\mathcal P(\omega)^M,\qquad
 \mathcal R_{\operatorname{HOD}}^M=
 \mathcal R^M\cap\operatorname{HOD}^M.
\]
The reported interface facts are
\[
 \operatorname{HOD}^M\models\mathrm{ZFC},\qquad
 \operatorname{Ord}^{\operatorname{HOD}^M}=\operatorname{Ord}^M,
 \qquad \mathcal R_{\operatorname{HOD}}^M
 \text{ is a set well-orderable in }M.
\]
The last statement means that a well-ordering relation on that set belongs to \(M\); it is not merely a well-order supplied by an external universe. In particular, if \(M\) has no well-ordering of \(\mathcal R^M\), then
\[
 \mathcal R_{\operatorname{HOD}}^M\subsetneq\mathcal R^M,
 \qquad \operatorname{HOD}^M\subsetneq M.
\]
The aid applies this under \(M\models\mathrm{ZF}+\mathrm{AD}\), in particular under \(\mathrm{ZF}+\mathrm{AD}^{+}\). Here \(\mathrm{AD}\) is the axiom that every two-player perfect-information length-\(\omega\) game on natural numbers is determined. The precise strengthened determinacy package denoted by \(\mathrm{AD}^{+}\) is not specified in the record; only its implication of \(\mathrm{AD}\) and the resulting non-well-orderability are used in this diagnostic.

Thus, for a proposed full presentation \(N\subseteq M\) claimed to equal \(\operatorname{HOD}^M\) in this determinacy setting, the aid identifies three necessary tests:
\begin{enumerate}
\item \(N\) must be an inner model satisfying \(\mathrm{ZFC}\) and containing all ordinals of \(M\).
\item Its ambient \(M\)-reals \(\mathcal R^M\cap N\) must be well-orderable by a relation in \(M\).
\item It must omit some real of \(M\).
\end{enumerate}
Failure of a test refutes the claimed full equality; passing all three is explicitly not sufficient. These elementary conditions do not construct an extender sequence, supply an iteration strategy, identify HOD, or solve a particular core-model problem. All ordinals, definability and real-set comparisons are relative to the fixed \(M\), including when \(M\) is a transitive set model rather than an ambient proper-class universe.
\subsection{Short English statement}
Yoshinobu's item is a statement of research interests. The supplied aid adds necessary checks for a proposed HOD presentation in a determinacy model: it must satisfy choice, contain the model's ordinals, have internally well-orderable reals, and omit an ambient real. These checks are not a sufficient characterization or an extender/strategy construction.
\subsection{Sources}
\begin{itemize}
\item Local UnsolvedMath registry, permanent ID \(20002384\), catalogue entry, \texttt{aim-logic-notes.json}, \texttt{notes[159]}, \emph{Recent Advances in Core Model Theory}, Chapter A, ``Participant Contributions'', item A.12, Yasuo Yoshinobu: \url{https://aimath.org/WWN/coremodel/coremodel.pdf}.
\item Supplied research summary attributes the organizers' interests to the matching official AIM announcement, including Woodin's HOD analysis in \(\mathrm{AD}^{+}\) models. No announcement URL, complete bibliographic reference or specific presentation theorem is supplied.
\item Supplied HOD diagnostic: attempt 1, context\_only, not\_a\_problem, rigor ``proved'', low novelty confidence. The separate open/OPEN-TRIAGE registry labels do not make the participant statement a posed theorem. The machine-generated research aid is not peer reviewed and must be independently verified before citation. Its source-verification field reports checking the official PDF; this draft has not independently visited it.
\end{itemize}
\end{document}
